\setlength{\tabcolsep}{4pt}
El Formulario T-8 identifica a las personas que LafroX asigna al contrato y distingue los servicios que ejecutará un tercero de los suministros con respaldo de fabricante. El Subdocumento 12, secciones 12.1 a 12.3, fundamenta ambas declaraciones.

\section{Equipo de trabajo}

Las ocho filas iniciales de la Tabla~\ref{tab:T8-equipo} corresponden al equipo clave; bajo «Otros roles» figuran los titulares funcional y de integración. La dedicación se lee en el orden implementación / Operación, como promedio de la fase. El guion indica que el rol no participa en esa fase.

\parte{tablas/tab_T8_equipo}

El Jefe de Proyecto tiene dedicación exclusiva en implementación. Para el Arquitecto de Solución se declara el promedio por arquitecto senior, con dedicación alta en diseño y asistencia permanente al Comité de Arquitectura. La fila «Líder de Calidad» corresponde al Líder de Calidad y Pruebas. Los Anexos 12.B a 12.D vinculan cada nominación con su currículo, sus horas y su compromiso individual.

\section{Subcontrataciones y alianzas}

LafroX contrata resultados especializados para la sala técnica y para la evaluación independiente de seguridad. La Tabla~\ref{tab:T8-servicios} distingue esos servicios de las cuatro alianzas de suministro y soporte; el porcentaje de subcontratación corresponde al conjunto, no a cada fila.

\parte{tablas/tab_T8_servicios}

LafroX seleccionará a cada empresa subcontratista mediante tres cotizaciones y solo adjudicará el servicio con autorización previa y escrita del CLIENTE. Los fabricantes actúan como último nivel de escalamiento: no asumen paquetes de la EDT ni sustituyen al responsable de LafroX ante el CLIENTE.
